% Generated by scripts/assemble.py; corrections belong in source/corrections.json.
\examhead{92}{南京师范大学 2025 年数学分析试卷}{91}
% p092-b01; source PDF page 92
\par\Needspace{5\baselineskip}
\blockanchor{p092-b01}
\majorquestion{一}综合计算题. 每题 5 分, 共 40 分.\par
% p092-b02; source PDF page 92
\blockanchor{p092-b02}
\question{1}$\lim_{x\to 1}\left(\frac{1}{x - 1} -\frac{1}{\ln x}\right).$\par
% p092-b03; source PDF page 92
\blockanchor{p092-b03}
\question{2}\displayformula{\lim _ {n \rightarrow \infty} n \left(\frac {1}{n ^ {2} + 1} + \frac {1}{n ^ {2} + 2 ^ {2}} + \dots + \frac {1}{n ^ {2} + n ^ {2}}\right).}\par
% p092-b04; source PDF page 92
\blockanchor{p092-b04}
\question{3}$\lim_{n\to\infty}\frac{n+1}{\sqrt[n]{n!}}$ .\par
% p092-b05; source PDF page 92
\blockanchor{p092-b05}
\question{4}设 $f(x) = \int_{0}^{x} \sin (x - t)e^{-t} \, \mathrm{d}t$, 求 $f'(x)$.\par
% p092-b06; source PDF page 92
\blockanchor{p092-b06}
\question{5}$\int e^{x}\cos^{2}x\mathrm{d}x.$\par
% p092-b07; source PDF page 92
\blockanchor{p092-b07}
\question{6}讨论 $\int_0^{+\infty} \frac{x^m}{1 + x^n} \, \mathrm{d}x$ 的敛散性, 其中 $m, n \in \mathbb{N}_+$.\par
% p092-b08; source PDF page 92
\blockanchor{p092-b08}
\question{7}设函数 $f$ 具有二阶连续偏导数, 令 $z = f(x + y, \arctan(xy))$, 求 $\frac{\partial^2 z}{\partial x \partial y}$.\par
% p092-b09; source PDF page 92
\blockanchor{p092-b09}
\question{8}计算曲面积分 $\iint_{\Sigma} (x + y + z)^2 \mathrm{~d} S$, 其中 $\Sigma: x^2 + y^2 + z^2 = R^2 (R > 0)$.\par
% p092-b10; source PDF page 92
\par\Needspace{5\baselineskip}
\blockanchor{p092-b10}
\majorquestion{二}每题 10 分, 共 20 分.\par
% p092-b11; source PDF page 92
\blockanchor{p092-b11}
\question{1}设 $a > 0, b > 0, a_1 = a, a_{n+1} = \frac{1}{2}\left(a_n + \frac{b}{a_n}\right), n \in \mathbb{N}_+$. 证明: 数列 $\{a_n\}$ 收敛并计算 $\lim_{n \to \infty} a_n$.\par
% p092-b12; source PDF page 92
\blockanchor{p092-b12}
\question{2}设函数 $f(x)$ 在 $[a, b]$ 上连续, 在 $(a, b)$ 上可导, 且 $f(a) = f(b) = 0$. 证明: 存在 $\xi \in (a, b)$, 使得\displayformula{f\,\prime(\xi)-2f(\xi)=0.}\par
% p092-b13; source PDF page 92
\par\Needspace{5\baselineskip}
\blockanchor{p092-b13}
\majorquestion{三}(15 分) 证明函数 $y = \sqrt{x}$ 在 $[0, +\infty)$ 上一致连续.\par
% p092-b14; source PDF page 92
\par\Needspace{5\baselineskip}
\blockanchor{p092-b14}
\majorquestion{四}(15 分) 求函数项级数 $\sum_{n=1}^{\infty} \frac{nx}{(1+x)(1+2x)\cdots(1+nx)}$ ($x \geq 0$) 的和函数, 并讨论其在 $[1, +\infty)$ 上的一致收敛性.\par
% p092-b15; source PDF page 92
\par\Needspace{5\baselineskip}
\par\Needspace{6\baselineskip}
\blockanchor{p092-b15}
\majorquestion{五}(15 分) 设 $f(u)$ 和 $g'(v)$ 在 $\mathbb{R}$ 上连续, $L: x^2 + y^2 = 4$, 取逆时针. $D: x^2 + y^2 \leq 4$, 且\par
\nopagebreak[4]
% p092-b16; source PDF page 92
\blockanchor{p092-b16}
\displayformula{\iint_ {D} (x + y) g ^ {\prime} (x - y) \mathrm{d} x \mathrm{d} y = 1.}
% p092-b17; source PDF page 92
\blockanchor{p092-b17}
计算 $I = \int_{L}[f(x^{2} + y^{2}) + g(x - y)](x\mathrm{d}x + y\mathrm{d}y).$\par
% p092-b18; source PDF page 92
\par\Needspace{5\baselineskip}
\par\Needspace{6\baselineskip}
\blockanchor{p092-b18}
\majorquestion{六}(15 分) 计算\par
\nopagebreak[4]
% p092-b19; source PDF page 92
\blockanchor{p092-b19}
\displayformula{I = \iint_ {\Sigma} 2 x ^ {3} \mathrm{d} y \mathrm{d} z + 2 y ^ {3} \mathrm{d} z \mathrm{d} x + 3 (z ^ {2} - 1) \mathrm{d} x \mathrm{d} y.}
% p092-b20; source PDF page 92
\blockanchor{p092-b20}
其中 $\Sigma$ 是曲面 $z = 1 - x^{2} - y^{2}$ ($z \geq 0$) 的上侧.\par
% p092-b21; source PDF page 92
\par\Needspace{5\baselineskip}
\blockanchor{p092-b21}
\majorquestion{七}(15 分) 设 $x_{n} \neq 0, n \in \mathbb{N}_{+}$, 证明: 若 $\lim_{n \to \infty} n\left(1 - \frac{|x_{n+1}|}{|x_{n}|}\right) = a > 1$, 则 $\sum_{n=1}^{\infty} x_{n}$ 绝对收敛.\par
% p092-b22; source PDF page 92
\par\Needspace{5\baselineskip}
\blockanchor{p092-b22}
\majorquestion{八}(15 分) 解答如下问题.\par
% p092-b23; source PDF page 92
\blockanchor{p092-b23}
(1) 叙述实数系完备性基本定理中的致密性定理和 Cauchy 收敛准则.\par
% p092-b24; source PDF page 92
\blockanchor{p092-b24}
(2) 利用致密性定理证明 Cauchy 收敛准则的充分性部分.\par
